% Bagean: sajarah
% Bab: biografi
% Pérangan: alfred-tarski

\documentclass[../../../include/open-logic-section]{subfiles}

\begin{document}

\olfileid{his}{bio}{tar}

\olsection{Alfred Tarski}

\olphoto{tarski-alfred}{Alfred Tarski}

Alfred Tarski lair tanggal 14 Januari 1901
ing Warsawa, Polandia (wektu iku bagéan
saka Kakaisaran Rusia). Tarski, kang
dicritakaké ``kaya Napoleon,'' iku
rame, seneng guneman, lan sregep banget.
Semangaté kerep katon ing kuliah-kuliahé---
dhèwèké tau ngobong kranjang sampah
nalika mbuwang rokok sajroning kuliah,
lan sawisé iku dilarang mulang ing
gedhong kasebut manèh.

Tarski wis ngelak kawruh wiwit enom.
Sanajan nalika wis tuwa dhèwèké
nyritakaké marang murid-muridé yèn
milih sinau logika amarga mung ing
kuliah iku dhèwèké olèh biji~B,
cathetan sekolah menengahé nuduhaké
yèn kabèh bijiné A---kalebu logika.
Dhèwèké kuliah ing Universitas
Warsawa wiwit taun 1918 nganti 1924.
Wiwitané Tarski arep sinau biologi,
nanging banjur kasengsem marang
matématika, filsafat, lan logika,
amarga universitas iku dadi pusat
Aliran Logika lan Filsafat Warsawa.
Tarski nggayuh gelar doktor
ing taun 1924 ing sangisoré
bimbingan Stanis\l{}aw
Le\'{s}niewski.

Sadurungé pindhah menyang Amérika
Sarékat ing taun 1939, Tarski
ngrampungaké sawetara karya
kang paling wigati nalika
nyambut gawé minangka guru
sekolah menengah ing Warsawa.
Karyané ngenani akibat logis
lan kabeneran logis ditulis
ing mangsa iku. Ing taun 1939
Tarski lunga menyang Amérika
Sarékat kanggo tur kuliah.
Nalika ana ing kono, Jerman
nyerang Polandia, lan amarga
asal-usulé Yahudi, Tarski
ora bisa bali. Garwané lan
anak-anaké tetep ana ing
Polandia nganti perang rampung,
nanging sawisé iku uga bisa
pindhah menyang Amérika Sarékat.
Tarski mulang ing Harvard,
College of the City of New York,
lan Institute for Advanced
Study ing Princeton, lan
pungkasané ing University
of California, Berkeley.
Ing kono dhèwèké ngedegaké
program lintas dhisiplin
Logika lan Metodologi Ilmu
Pengetahuan. Tarski séda
tanggal 26 Oktober 1983
nalika umur 82 taun.

\begin{reading}
Kanggo katrangan luwih jangkep
ngenani uripé Tarski, delengen
biografi \emph{Alfred Tarski: Life
  and Logic} \citep{Feferman2004}.
Karya-karya utama Tarski ngenani
akibat logis lan kabeneran
kasedhiya ing basa Inggris
ing \citep{Tarski1983}.
Kabèh karya asli Tarski
wis diklumpukaké dadi
seri patang jilid,
\citep{Tarski1981}.
\end{reading}

\end{document}
